% 3.模型假设.tex —— 三、模型假设
\section{模型假设}
为简化问题，本文做出以下假设：
\begin{itemize}
  \item \textbf{假设 1（几何）}：药材为均质、各向同性的无限长圆柱，温度与水分浓度仅沿径向变化。
  \item \textbf{假设 2（物性）}：药材密度 $\rho$、比热容 $c_p$、导热系数 $k$ 在所考察温度与含水率范围内取常值。
  \item \textbf{假设 3（初始状态）}：药材内部初始温度与含水率均匀，分别为 $T_0=28$ $^\circ$C、$C_0=2.55$ kg/kg（干基）。
  \item \textbf{假设 4（边界）}：药材表面与烘房环境之间的换热与传质均服从牛顿冷却定律形式，即通量与表面—环境之差的一次方成正比。
  \item \textbf{假设 5（无内源）}：药材内部无相变潜热释放、无化学放热、无水分生成。这是本问最主要的简化——干燥中水分汽化所需潜热由药材自身供给，因题设未给出汽化潜热，该项被略去。
  \item \textbf{假设 6（扩散形式）}：水分迁移服从 Fick 定律，有效扩散系数 $D$ 仅依赖局部含水率 $C$。
\end{itemize}
